
\subsection{$\mathcal M_{k+1}^{\mathrm{main,reg}}(\beta)$
之外的仓西坐标图：胶合}
\label{subsec:7-1-3}

上一小节已经为
$\mathcal M_{k+1}^{\mathrm{main,reg}}(\beta)$
的每一点构造了仓西邻域。
现在要为
$\mathcal M_{k+1}^{\mathrm{main}}(\beta)
\setminus\mathcal M_{k+1}^{\mathrm{main,reg}}(\beta)$
中的点构造仓西邻域；为此需要进行陶布斯型胶合构造\liftIndex{I229}。

首先定义图的胶合映射
\[
\begin{aligned}
\operatorname{Comp}_{k_1\cdots k_m;m';j_1,\ldots,j_m}\colon
&\overline{\operatorname{Gr}}_{k_1+1}\times\cdots\times
 \overline{\operatorname{Gr}}_{k_m+1}\times
 \overline{\operatorname{Gr}}_{m+m'+1}\\
&\longrightarrow
 \overline{\operatorname{Gr}}_{k_1+\cdots+k_m+m'+1},
\end{aligned}
\]
其中 $0<j_1<j_2<\cdots<j_m<m+m'+1$。
设 $\mathfrak t_i\in G_{k_i+1}$（$i=1,\ldots,m$），
$\mathfrak t_0\in G_{m+m'+1}$。
把 $\mathfrak t_i$ 的第 $0$ 条外边
与 $\mathfrak t_0$ 的第 $j_i$ 条边相识别，
得到 $\mathfrak t\in G_{\sum_i k_i+m'+1}$。
给定 $\ell_i\in\operatorname{Gr}(\mathfrak t_i)$
及 $\ell_0\in\operatorname{Gr}(\mathfrak t_0)$，
如下定义映射
$\ell\colon C^1_{\mathrm{int}}(\mathfrak t)\to(0,\infty]$：
\[
C^1_{\mathrm{int}}(\mathfrak t)
=\bigcup_i C^1_{\mathrm{int}}(\mathfrak t_i)
 \cup C^1_{\mathrm{int}}(\mathfrak t_0)
 \cup\{e_1,\ldots,e_m\},
\]
其中 $e_i$ 是由 $\mathfrak t_i$ 的第 $0$ 条外边产生的新边，并令
\[
\ell(e)=
\begin{cases}
\ell_i(e),&e\in C^1_{\mathrm{int}}(\mathfrak t_i),\\
\ell_0(e),&e\in C^1_{\mathrm{int}}(\mathfrak t_0),\\
\infty,&e=e_i,\quad i=1,\ldots,m.
\end{cases}
\]
定义
\[
\operatorname{Comp}_{k_1\cdots k_m;m';j_1,\ldots,j_m}
\bigl((\mathfrak t_1,\ell_1),\ldots,
(\mathfrak t_m,\ell_m);(\mathfrak t_0,\ell_0)\bigr)
=:(\mathfrak t,\ell).
\]
每个这样的 $(\mathfrak t,\ell)$
都能唯一地由模空间内部的
$\operatorname{Gr}_{*}$ 元素反复施用
$\operatorname{Comp}$ 得到；见图~\ref{fig:7-1-2}。
本小节余下部分只考虑
\[
(\mathfrak t,\ell)
=\operatorname{Comp}_{k;m;j}
  \bigl((\mathfrak t_1,\ell_1),(\mathfrak t_0,\ell_0)\bigr),
\quad
(\mathfrak t_1,\ell_1)\in\operatorname{Gr}_{k+1},\quad
(\mathfrak t_0,\ell_0)\in\operatorname{Gr}_{m+1}
\]
的情形。

\par\medskip\noindent\begin{minipage}{\textwidth}
\centering
\setlength{\unitlength}{0.70mm}
\begin{picture}(145,56)
\linethickness{0.48pt}
\put(3,42){\line(1,0){15}}
\put(18,42){\line(2,1){16}}
\put(18,42){\line(2,-1){16}}
\put(18,42){\line(0,-1){18}}
\put(34,34){\line(1,0){13}}
\put(19,48){$\scriptstyle\mathfrak t_0$}
\put(36,26){$\scriptstyle j$}
\put(3,9){\line(1,0){15}}
\put(18,9){\line(2,1){16}}
\put(18,9){\line(2,-1){16}}
\put(18,9){\line(0,1){11}}
\put(6,20){$\scriptstyle\mathfrak t_1$}
\put(58,27){$\Longrightarrow$}
\put(78,42){\line(1,0){15}}
\put(93,42){\line(2,1){16}}
\put(93,42){\line(2,-1){16}}
\put(93,42){\line(0,-1){18}}
\put(109,34){\line(1,0){15}}
\put(93,24){\line(2,-1){17}}
\put(110,15){\line(2,1){17}}
\put(110,15){\line(2,-1){17}}
\put(110,15){\line(0,1){9}}
\put(89,10){$\scriptstyle e(\mathrm{new})$}
\end{picture}
\par\smallskip
{\renewcommand{\thefigure}{7.1.2}
\refstepcounter{figure}
{\normalfont\bfseries 图~\thefigure}\label{fig:7-1-2}}
\end{minipage}\par\medskip

换言之，在映射
$\Theta\colon\overline{\operatorname{Gr}}_{k+m}
\to\mathcal M^{\mathrm{main}}_{k+m}$ 下，
$(\mathfrak t,\ell)$ 对应一条亏格为 $0$、
边界上恰有一个奇点（即一个节点）的
带标记点的半稳定曲线。
\translatornote{原书此处把本专门胶合情形的 $\Theta$ 型别写成下标 $k+1$，后文又把总模空间写成 $\mathcal M^{\mathrm{main}}_{k+m+1}$；两分支各消耗一条外边，故总标记数应为 $k+m$，也与随后 $\operatorname{Gr}_{k+m}$ 相合。见疑误 A-0413。}
其他情形的处理相同，只是记号更复杂。
球面气泡也可依照 \textbf{[FuOn99II]} 的方法处理，
参见第~7.1.4 小节。

设
$w\colon(X(\mathfrak t,\ell),\partial X(\mathfrak t,\ell))
\to(M,L)$ 是伪全纯映射；
当 $X(\mathfrak t,\ell)$ 双全纯于
$\Sigma\setminus\{\text{标记点}\}$ 时，
它延拓成稳定映射
$(\Sigma,\partial\Sigma)\to(M,L)$，
见第~3.4 节。
于是 $w$ 确定映射
\[
w_1\colon X(\mathfrak t_1,\ell_1)\to M,
\qquad
w_0\colon X(\mathfrak t_0,\ell_0)\to M,
\]
以及点
\[
\vec p^{\,1}=(p_0^1,\ldots,p_k^1)\in L^{k+1},
\qquad
\vec p^{\,0}=(p_0^0,\ldots,p_m^0)\in L^{m+1},
\]
使得
\[
\begin{gathered}
((\mathfrak t_1,\ell_1),w_1,\vec p^{\,1})
 \in W_{k+1}^{1,p}(M,L),\\
((\mathfrak t_0,\ell_0),w_0,\vec p^{\,0})
 \in W_{m+1}^{1,p}(M,L),
\qquad p_0^1=p_j^0.
\end{gathered}
\]
如第~7.1.2 小节，
我们已选取
$E_{((\mathfrak t_1,\ell_1),w_1,\vec p^{\,1})}$
与
$E_{((\mathfrak t_0,\ell_0),w_0,\vec p^{\,0})}$。
定义如下空间的子空间 $\mathcal X$：
\begin{equation}
\begin{aligned}
\mathcal X^+:={}&
T_{((\mathfrak t_1,\ell_1),w_1,\vec p^{\,1})}
 W_{k+1}^{1,p}(M,L)\\
&\oplus
T_{((\mathfrak t_0,\ell_0),w_0,\vec p^{\,0})}
 W_{m+1}^{1,p}(M,L).
\end{aligned}
\tag{7.1.19.1}\label{eq:7-1-19-1}
\end{equation}
具体地，
\begin{equation}
\begin{aligned}
\mathcal X:=\bigl\{&((\xi_1,V_1,\vec v^{\,1}),
 (\xi_0,V_0,\vec v^{\,0}))\in\mathcal X^+\;\big|\\
&v_0^1=v_j^0\in T_{p_0^1}L=T_{p_j^0}L\bigr\}.
\end{aligned}
\tag{7.1.19.2}\label{eq:7-1-19-2}
\end{equation}
这里第二分量取 $V_0$。
\translatornote{原书式 (7.1.19.2) 及本页后段均将第二分量的 $V_0$ 写成 $V_1$；第二直和项以 $w_0$ 为基点，故依其型别改正。见疑误 A-0410。}
将
$D_{((\mathfrak t_1,\ell_1),w_1,\vec p^{\,1})}\bar\partial$
与
$D_{((\mathfrak t_0,\ell_0),w_0,\vec p^{\,0})}\bar\partial$
的直和限制到 $\mathcal X$，记为
\[
\begin{aligned}
D_{(w_1,w_0)}\bar\partial\colon\mathcal X
\longrightarrow{}&
\mathcal E^{0,p}_{((\mathfrak t_1,\ell_1),w_1,\vec p^{\,1})}\\
&\oplus
\mathcal E^{0,p}_{((\mathfrak t_0,\ell_0),w_0,\vec p^{\,0})}.
\end{aligned}
\]

\begin{lemma}{7.1.20}\label{lem:7-1-20}
复合映射
\[
\begin{aligned}
\pi\circ D_{(w_1,w_0)}\bar\partial\colon\mathcal X
\longrightarrow{}&
\mathcal E^{0,p}_{((\mathfrak t_1,\ell_1),w_1,\vec p^{\,1})}
 /E_{((\mathfrak t_1,\ell_1),w_1,\vec p^{\,1})}\\
&\oplus
\mathcal E^{0,p}_{((\mathfrak t_0,\ell_0),w_0,\vec p^{\,0})}
 /E_{((\mathfrak t_0,\ell_0),w_0,\vec p^{\,0})}
\end{aligned}
\]
是满射。其中 $\pi$ 是自然投影。
\end{lemma}

\begin{proof}
这由引理~\ref{lem:7-1-18}
和姆罗夫卡的 ``Mayer--Vietoris 原理'' \textbf{[Mro89]}
推出；亦可与 \textbf{[Fuk96II]} 的图~6 比较。
\end{proof}

利用引理~\ref{lem:7-1-20}，
现在进行陶布斯型胶合构造。
须把伪全纯曲线方程改写成类似式~\eqref{eq:7-1-16}
的方程，并构造其解的模空间。
此模空间由 $\pi\circ D_u\bar\partial$ 的核，
连同 $(\mathfrak t,\ell)$ 在
$\mathcal M^{\mathrm{main}}_{k+m}$
中的一个邻域参数化。
后一邻域由 $\mathcal V\times(C,\infty]$ 参数化，
其中 $\mathcal V$ 是
\[
T_{(\mathfrak t_1,\ell_1)}
\mathcal M^{\mathrm{main}}_{k+1}
\oplus
T_{(\mathfrak t_0,\ell_0)}
\mathcal M^{\mathrm{main}}_{m+1}
\]
中 $0$ 的一个邻域，
$(C,\infty]$ 则表示控制边界奇点处
相接圆盘的“拉伸”的参数集合；
$C$ 取充分大的正数。

如同其他胶合构造，我们先用单位分解
构造稍后定义的式 (7.1.22) 的近似解，
再研究其线性化并检验满射性。
由于在无穷远处发生退化，
要使所需估计一致成立，
须采用带有某种非通常权重的加权索伯列夫范数。
这一方法类似 \textbf{[Fuk96II]} 第~8 节所用的方法。

令
$(\mathfrak t'_1,\ell'_1)\in\operatorname{Gr}_{k+1}$、
$(\mathfrak t'_0,\ell'_0)\in\operatorname{Gr}_{m+1}$
分别位于
$(\mathfrak t_1,\ell_1)$、
$(\mathfrak t_0,\ell_0)$ 的小邻域中。
对于每个充分大的正数 $T$，
以及
$\ker(\pi\circ D_u\bar\partial)$ 中的元
$((\xi_1,V_1,\vec v^{\,1}),
(\xi_0,V_0,\vec v^{\,0}))$，
我们要构造伪全纯曲线方程的近似解。

先胶合映射的定义域，取
\[
(\mathfrak t',\ell'_\infty)
=\operatorname{Comp}_{k;m;j}
 \bigl((\mathfrak t'_1,\ell'_1),
       (\mathfrak t'_0,\ell'_0)\bigr).
\]
图 $(\mathfrak t',\ell'_\infty)$ 中有一条无限长的内边，
即胶合 $\mathfrak t'_0$ 与 $\mathfrak t'_1$ 之处的边；
记为 $e(\mathrm{new})$。
把它替换成长度为 $10T$ 的边，所得图记为
\[
(\mathfrak t',\ell'_T)
=\operatorname{Comp}
 \bigl((\mathfrak t'_1,\ell'_1),
       (\mathfrak t'_0,\ell'_0);T\bigr)
\in\operatorname{Gr}_{k+m}.
\]
更确切地，
$X(\mathfrak t',\ell'_T)$ 由
$X(\mathfrak t'_1,\ell'_1)$ 和
$X(\mathfrak t'_0,\ell'_0)$ 如下构成。
以 $e_{1,0}$ 表示 $\mathfrak t'_1$ 的第 $0$ 条外边，
以 $e_{0,j}$ 表示 $\mathfrak t'_0$ 的第 $j$ 条外边。
为便于说明，重定这两个端的参数，使得
\[
\begin{aligned}
L_{e_{1,0}}&\cong[-5T,\infty)\times[0,1]
 \subset X(\mathfrak t'_1,\ell'_1),\\
L_{e_{0,j}}&\cong(-\infty,5T]\times[0,1]
 \subset X(\mathfrak t'_0,\ell'_0),
\end{aligned}
\]
而不再用原参数化
$L_{e_{1,0}}\cong[0,\infty)\times[0,1]$、
$L_{e_{0,j}}\cong(-\infty,0]\times[0,1]$。
从 $X(\mathfrak t'_1,\ell'_1)$ 中删去
$[0,\infty)\times[0,1]\subset L_{e_{1,0}}$，
从 $X(\mathfrak t'_0,\ell'_0)$ 中删去
$(-\infty,0]\times[0,1]\subset L_{e_{0,j}}$，
再沿 $\{0\}\times[0,1]$ 胶合。
这时
\[
L_{e(\mathrm{new})}\cong[-5T,5T]\times[0,1]
\subset X(\mathfrak t',\ell'_T),
\]
故所得正是 $X(\mathfrak t',\ell'_T)$。
\translatornote{原书在此及图 7.1.3 把胶合后域误写成带 $\mathfrak t'_0$ 的域；随后截断函数与截面段还混用 $\ell'$ 或无撇的 $(\mathfrak t,\ell)$。这里依紧前定义与自然映射的目标统一为 $X(\mathfrak t',\ell'_T)$。见疑误 A-0411。}

\par\medskip\noindent\begin{minipage}{\textwidth}
\centering
\setlength{\unitlength}{0.70mm}
\begin{picture}(145,61)
\linethickness{0.46pt}
\put(7,53){$\scriptstyle X(\mathfrak t'_1,\ell'_1)$}
\put(7,43){\line(2,1){8}}
\put(7,43){\line(2,-1){8}}
\put(15,41){\line(1,0){49}}
\put(15,45){\line(1,0){49}}
\put(30,40){\line(0,1){6}}
\put(24,34){$\scriptstyle-5T$}
\put(78,53){$\scriptstyle X(\mathfrak t'_0,\ell'_0)$}
\put(77,41){\line(1,0){49}}
\put(77,45){\line(1,0){49}}
\put(126,43){\line(2,1){8}}
\put(126,43){\line(2,-1){8}}
\put(113,40){\line(0,1){6}}
\put(108,34){$\scriptstyle5T$}
\put(70,26){$\Downarrow$}
\put(7,13){\line(2,1){8}}
\put(7,13){\line(2,-1){8}}
\put(15,11){\line(1,0){111}}
\put(15,15){\line(1,0){111}}
\put(126,13){\line(2,1){8}}
\put(126,13){\line(2,-1){8}}
\put(30,10){\line(0,1){6}}
\put(113,10){\line(0,1){6}}
\put(25,17){$\scriptstyle-5T$}
\put(108,17){$\scriptstyle5T$}
\put(58,12){$\scriptstyle L_{e(\mathrm{new})}$}
\put(50,2){$\scriptstyle X(\mathfrak t',\ell'_T)$}
\end{picture}
\par\smallskip
{\renewcommand{\thefigure}{7.1.3}
\refstepcounter{figure}
{\normalfont\bfseries 图~\thefigure}\label{fig:7-1-3}}
\end{minipage}\par\medskip

由构造可得自然映射
\[
\begin{aligned}
I_{(\mathfrak t_0,\ell_0),(\mathfrak t',\ell'_T)}\colon
&X(\mathfrak t_0,\ell_0)
 \setminus\bigl((-\infty,-5T)\times[0,1]\bigr)
 \longrightarrow X(\mathfrak t',\ell'_T),\\
I_{(\mathfrak t_1,\ell_1),(\mathfrak t',\ell'_T)}\colon
&X(\mathfrak t_1,\ell_1)
 \setminus\bigl((5T,\infty)\times[0,1]\bigr)
 \longrightarrow X(\mathfrak t',\ell'_T).
\end{aligned}
\]
其中 $(5T,\infty)\times[0,1]$
视为 $L_{e_{1,0}}$ 的子集，
$(-\infty,-5T)\times[0,1]$
视为 $L_{e_{0,j}}$ 的子集。
显然，这两个 $I$ 映射各自到其像都是微分同胚；
两像之并为 $X(\mathfrak t',\ell'_T)$，
交集为 $L_{e(\mathrm{new})}$，
见图~\ref{fig:7-1-3}。
接下来引入若干截断函数，
见图~\ref{fig:7-1-4}。

\par\medskip\noindent\begin{minipage}{\textwidth}
\centering
\setlength{\unitlength}{0.69mm}
\begin{picture}(145,73)
\linethickness{0.45pt}
\put(8,59){\vector(1,0){57}}
\put(14,55){\vector(0,1){17}}
\put(17,59){\line(1,0){21}}
\qbezier(38,59)(43,59)(45,69)
\qbezier(45,69)(47,72)(51,72)
\put(51,72){\line(1,0){11}}
\put(11,71){$\scriptstyle1$}
\put(29,54){$\scriptstyle S-1$}
\put(43,54){$\scriptstyle S$}
\put(50,54){$\scriptstyle S+1$}
\put(53,65){$\scriptstyle\chi_S^+$}
\put(74,59){\vector(1,0){61}}
\put(80,55){\vector(0,1){17}}
\put(82,72){\line(1,0){23}}
\qbezier(105,72)(110,72)(112,63)
\qbezier(112,63)(114,59)(118,59)
\put(118,59){\line(1,0){14}}
\put(77,71){$\scriptstyle1$}
\put(98,54){$\scriptstyle S-1$}
\put(110,54){$\scriptstyle S$}
\put(118,54){$\scriptstyle S+1$}
\put(119,66){$\scriptstyle\chi_S^-$}
\put(10,42){\line(1,0){52}}
\put(10,35){\line(1,0){52}}
\put(18,44){$\scriptstyle X(\mathfrak t_1,\ell_1)$}
\put(20,37){$\scriptstyle\chi_S^1\!\approx1$}
\put(54,37){$\scriptstyle0$}
\put(76,42){\line(1,0){59}}
\put(76,35){\line(1,0){59}}
\put(91,44){$\scriptstyle X(\mathfrak t_0,\ell_0)$}
\put(78,37){$\scriptstyle0$}
\put(115,37){$\scriptstyle\chi_S^0\!\approx1$}
\put(29,21){\line(1,0){86}}
\put(29,14){\line(1,0){86}}
\put(49,23){$\scriptstyle X(\mathfrak t',\ell'_T)$}
\put(35,16){$\scriptstyle\chi_S^-\!\approx1$}
\put(92,16){$\scriptstyle\chi_S^+\!\approx1$}
\put(69,14){\line(0,1){7}}
\put(70,8){$\scriptstyle S$}
\end{picture}
\par\smallskip
{\renewcommand{\thefigure}{7.1.4}
\refstepcounter{figure}
{\normalfont\bfseries 图~\thefigure}\label{fig:7-1-4}}
\end{minipage}\par\medskip

对每个 $S\in[-5T+1,5T-1]$，
取光滑函数
$\chi_S^+\colon[-5T,5T]\to[0,1]$，使
\[
\chi_S^+(\tau)=
\begin{cases}
0,&\tau<S-1,\\
1,&\tau>S+1,
\end{cases}
\]
并满足
\begin{equation}
|\nabla\chi_S^+|<C.
\tag{7.1.21}\label{eq:7-1-21}
\end{equation}
令 $\chi_S^-=1-\chi_S^+$。
把它们看成定义在
$L_{e(\mathrm{new})}=[-5T,5T]\times[0,1]$
上的函数后，
分别在
$X(\mathfrak t',\ell'_T)\setminus L_{e(\mathrm{new})}$
的两个分支上取 $0$ 或 $1$，
便将 $\chi_S^\pm$ 延拓到整个
$X(\mathfrak t',\ell'_T)$；
延拓后仍用相同记号。注意
\[
\begin{aligned}
\chi_S^1&=\chi_S^-\circ
I_{(\mathfrak t_1,\ell_1),(\mathfrak t',\ell'_T)},\\
\chi_S^0&=\chi_S^+\circ
I_{(\mathfrak t_0,\ell_0),(\mathfrak t',\ell'_T)}.
\end{aligned}
\]
另一方面，将 $\chi_S^-$ 看成
$L_{e_{1,0}}\cong[-5T,\infty)\times[0,1]$
上的函数，再以同样方式延拓到
$X(\mathfrak t_1,\ell_1)$，
所得函数记作 $\chi_S^1$；
将 $\chi_S^+$ 看成
$L_{e_{0,j}}\cong(-\infty,5T]\times[0,1]$
上的函数并延拓到 $X(\mathfrak t_0,\ell_0)$，
所得函数记作 $\chi_S^0$。

我们将用这些截断函数把
$X(\mathfrak t_0,\ell_0)$ 与
$X(\mathfrak t_1,\ell_1)$ 上的截面胶合成
$X(\mathfrak t',\ell'_T)$ 上的截面。
由于
$w_0\in\mathcal M^{\mathrm{main}}_{m+1}(\beta_0)$
与
$w_1\in\mathcal M^{\mathrm{main}}_{k+1}(\beta_1)$
在无穷远处指数收敛，
可在 $p_0^1=p_j^0$ 的一个小坐标邻域中
胶合二者，得到映射
$w'_T\colon X(\mathfrak t',\ell'_T)\to M$。
\translatornote{原书此处把 $w_0$、$w_1$ 所属模空间的标记数 $k+1$ 与 $m+1$ 写反；译文依前文两个分支各自的定义域改正。见疑误 A-0412。}

为写出接下来的四个映射，记
\[
\begin{aligned}
\mathcal F_i&:=\Gamma\bigl(X(\mathfrak t_i,\ell_i),
 \Lambda^{0,1}\otimes w_i^*TM\bigr),\quad i=0,1,\\
\mathcal F_T&:=\Gamma\bigl(X(\mathfrak t',\ell'_T),
 \Lambda^{0,1}\otimes (w'_T)^*TM\bigr).
\end{aligned}
\]
用于截面胶合的四个映射是
\[
\begin{array}{ll}
J_{0,*}^{S}\colon\mathcal F_0\longrightarrow\mathcal F_T,
&J_{1,*}^{S}\colon\mathcal F_1\longrightarrow\mathcal F_T,\\[2pt]
J_{*,0}^{S}\colon\mathcal F_T\longrightarrow\mathcal F_0,
&J_{*,1}^{S}\colon\mathcal F_T\longrightarrow\mathcal F_1.
\end{array}
\]
在 $I_{(\mathfrak t_0,\ell_0),(\mathfrak t',\ell'_T)}$
的像上，定义
\[
w'_0:=w_0\circ
I_{(\mathfrak t_0,\ell_0),(\mathfrak t',\ell'_T)}^{-1},
\]
并以同样方式在
$I_{(\mathfrak t_1,\ell_1),(\mathfrak t',\ell'_T)}$
的像上定义 $w'_1$。

设
$V\in\Gamma\bigl(X(\mathfrak t_0,\ell_0),
\Lambda^{0,1}\otimes w_0^*TM\bigr)$。
$\chi_S^0V$ 的支集包含于
\[
X(\mathfrak t_0,\ell_0)
\setminus\bigl((-\infty,-5T)\times[0,1]\bigr).
\]
因此可借助
$I_{(\mathfrak t_0,\ell_0),(\mathfrak t',\ell'_T)}$
把它送到胶合后域的相应像内，
成为沿 $w'_0$ 的截面；
再沿最短测地线作平行移动，
便得到 $J_{0,*}^{S}(V)\in\mathcal F_T$。
以 $\chi_S^1$ 代替 $\chi_S^0$
可同样定义 $J_{1,*}^{S}$。

反过来，设 $V\in\mathcal F_T$。
$\chi_S^+V$ 的支集位于
$I_{(\mathfrak t_0,\ell_0),(\mathfrak t',\ell'_T)}$
的像内。
沿最短测地线平行移动，
先将其变成沿 $w'_0$ 的截面；
再由该 $I$ 映射拉回，
得到 $J_{*,0}^{S}(V)\in\mathcal F_0$。
改用 $\chi_S^-$ 即可同样定义
$J_{*,1}^{S}$。
这些映射还延拓到相关截面空间的
$W^{1,p}$ 或 $W^{0,p}$ 闭包。
